\section*{अध्याय \ref{ch:g7:chemical-reactions} --- रासायनिक अभिक्रियाएँ: अभिकारक और उत्पाद}

\begin{solution}{exo:g7:chemical-reactions:1}
भौतिक: बर्फ़ का पिघलना, चीनी का \omterm{def:g2:mixing-and-dissolving:dissolve}{घुलना}। रासायनिक: तीली का \omterm{def:g4:what-a-fire-needs:burning}{जलना},
डबलरोटी का सिंकना, कील में जंग लगना।
\end{solution}

\begin{solution}{exo:g7:chemical-reactions:2}
\omterm{def:g7:chemical-reactions:reactant}{अभिकारक} वे स्पीशीज़ हैं जिनकी खपत होती है; \omterm{def:g7:chemical-reactions:reactant}{उत्पाद} वे नई स्पीशीज़ हैं जो
बनती हैं।
\end{solution}

\begin{solution}{exo:g7:chemical-reactions:3}
कार्बन $+$ ऑक्सीजन $\longrightarrow$ कार्बन डाइऑक्साइड।
\end{solution}

\begin{solution}{exo:g7:chemical-reactions:4}
\omterm{def:g7:chemical-reactions:reactant}{अभिकारक}: लोहा और ऑक्सीजन। \omterm{def:g7:chemical-reactions:reactant}{उत्पाद}: आयरन ऑक्साइड।
\end{solution}

\begin{solution}{exo:g7:chemical-reactions:5}
धुँधलापन पानी दिखाता है; दूधिया चूने का पानी कार्बन डाइऑक्साइड दिखाता है।
\end{solution}

\begin{solution}{exo:g7:chemical-reactions:6}
ज़िंक $+$ हाइड्रोक्लोरिक अम्ल $\longrightarrow$ ज़िंक क्लोराइड $+$ हाइड्रोजन।
\end{solution}

\begin{solution}{exo:g7:chemical-reactions:7}
पहले: 4 (मेथेन के \omterm{def:g7:atoms-and-molecules:molecule}{अणु} में)। बाद में: $2 + 2 = 4$ (पानी के हर \omterm{def:g7:atoms-and-molecules:molecule}{अणु}
में दो)।
\end{solution}

\begin{solution}{exo:g7:chemical-reactions:8}
कोई नई स्पीशीज़ प्रकट नहीं होती: नमक अब भी पानी में फैला हुआ मौजूद है,
और पानी का \omterm{def:g3:separating-mixtures:evaporate}{वाष्पन} करके वापस मिल जाता है। यह एक \omterm{def:g7:chemical-reactions:physical-transformation}{भौतिक
रूपांतरण} है।
\end{solution}

\begin{solution}{exo:g7:chemical-reactions:9}
ज़्यादातर लकड़ी ने ऑक्सीजन के साथ अभिक्रिया करके ऐसे \omterm{def:g7:chemical-reactions:reactant}{उत्पाद} दिए हैं जो गैसें
हैं, कार्बन डाइऑक्साइड और जलवाष्प, और वे \omterm{def:g4:air-a-mixture-of-gases:air}{वायु} में चले जाते हैं। केवल
राख बचती है।
\end{solution}

\begin{solution}{exo:g7:chemical-reactions:10}
हाइड्रोजन $+$ ऑक्सीजन $\longrightarrow$ पानी। \omterm{def:g7:chemical-reactions:reactant}{उत्पाद} निर्जल
कॉपर सल्फ़ेट को नीला कर देता है।
\end{solution}

\begin{solution}{exo:g7:chemical-reactions:11}
पहले: $2 \times 2 = 4$ हाइड्रोजन \omterm{def:g7:atoms-and-molecules:atom}{परमाणु} और 2 ऑक्सीजन \omterm{def:g7:atoms-and-molecules:atom}{परमाणु}। बाद में: पानी
के दो \omterm{def:g7:atoms-and-molecules:molecule}{अणुओं} में $2 \times 2 = 4$ हाइड्रोजन \omterm{def:g7:atoms-and-molecules:atom}{परमाणु} और $2 \times 1 =
2$ ऑक्सीजन \omterm{def:g7:atoms-and-molecules:atom}{परमाणु} हैं। वही \omterm{def:g7:atoms-and-molecules:atom}{परमाणु}, उतनी ही संख्या में।
\end{solution}

\begin{solution}{exo:g7:chemical-reactions:12}
कोई ऑक्सीजन \omterm{def:g7:atoms-and-molecules:atom}{परमाणु} कहीं से नहीं आ सकता: चित्रण में पहले 2 ऑक्सीजन \omterm{def:g7:atoms-and-molecules:atom}{परमाणु}
हैं और बाद में 3। सही: एक कार्बन \omterm{def:g7:atoms-and-molecules:atom}{परमाणु} $+$ ऑक्सीजन का एक \omterm{def:g7:atoms-and-molecules:molecule}{अणु}
$\longrightarrow$ कार्बन डाइऑक्साइड का एक \omterm{def:g7:atoms-and-molecules:molecule}{अणु}, और कुछ भी नहीं बचता।
\end{solution}

\begin{solution}{pb:g7:chemical-reactions:1}
\textbf{1.} मेथेन और ऑक्सीजन।

\textbf{2.} पानी।

\textbf{3.} कार्बन डाइऑक्साइड।

\textbf{4.} नहीं: उसकी खपत नहीं होती, और वह \omterm{def:g7:chemical-reactions:word-equation}{शाब्दिक समीकरण} में नहीं
लिखी जाती।

\textbf{5.} मेथेन $+$ ऑक्सीजन $\longrightarrow$ कार्बन डाइऑक्साइड $+$
पानी।

\textbf{6.} रासायनिक: दो स्पीशीज़ ग़ायब होती हैं (मेथेन, ऑक्सीजन), और परीक्षणों
से पहचानी गई दो नई स्पीशीज़ प्रकट होती हैं; और ऊष्मा तथा प्रकाश निकलते हैं।

\textbf{7.} \ce{CH4}, \ce{O2}, \ce{CO2}, \ce{H2O}।

\textbf{8.} 1 कार्बन, 4 हाइड्रोजन, $2 \times 2 = 4$ ऑक्सीजन।

\textbf{9.} 1 कार्बन; $2 \times 2 = 4$ हाइड्रोजन; $2 + 2 \times 1 = 4$
ऑक्सीजन। वही \omterm{def:g7:atoms-and-molecules:atom}{परमाणु}, उतनी ही संख्या में: वे केवल नए ढंग से
व्यवस्थित हुए हैं।

\textbf{10.} $10 \times 2 = 20$ ऑक्सीजन \omterm{def:g7:atoms-and-molecules:molecule}{अणु}।

\textbf{11.} कार्बन डाइऑक्साइड के 10 \omterm{def:g7:atoms-and-molecules:molecule}{अणु}।

\textbf{12.} $10 \times 2 = 20$ पानी के \omterm{def:g7:atoms-and-molecules:molecule}{अणु}।
\end{solution}
